%% ============================================================
%%  SECCIÓN 1 — Introducción
%%  Responsable: Vera Vivanco, Jesús Rodolfo
%% ============================================================

\section{Introducción}

Los materiales sólidos expuestos a variaciones de temperatura experimentan cambios en sus dimensiones geométricas. A nivel molecular, al incrementarse la temperatura de un sólido, se eleva la energía cinética de sus átomos, lo que provoca vibraciones con mayor amplitud alrededor de sus posiciones de equilibrio. Debido a la naturaleza asimétrica (anharmónica) del potencial de interacción interatómico, la distancia promedio entre los átomos vecinos aumenta a medida que aumenta la temperatura. A nivel macroscópico, este incremento acumulativo del espaciamiento atómico se manifiesta como una dilatación térmica de las dimensiones del cuerpo \cite{serway}.

Cuando este fenómeno se restringe predominantemente a una sola dimensión, se define como dilatación lineal. La magnitud física que caracteriza cuantitativamente este comportamiento es el \textit{coeficiente de dilatación lineal} $\alpha$, definido por la relación $\alpha = \frac{\Delta L}{L_0\,\Delta T}$, donde $\Delta L$ representa la variación de la longitud del cuerpo, $L_0$ su longitud inicial a la temperatura de referencia y $\Delta T$ la variación de temperatura aplicada. En el Sistema Internacional de Unidades (SI), el coeficiente $\alpha$ se expresa en unidades de \unit{\per\kelvin}, aunque es habitual utilizar la unidad equivalente de $^\circ\text{C}^{-1}$ \cite{halliday}. Para los materiales bajo estudio (aluminio, cobre y vidrio), los valores de referencia bibliográficos aceptados son $\alpha_{\text{ref, Al}} = 2{,}30 \times 10^{-5}\ ^\circ\text{C}^{-1}$, $\alpha_{\text{ref, Cu}} = 1{,}70 \times 10^{-5}\ ^\circ\text{C}^{-1}$ y $\alpha_{\text{ref, vidrio}} = 2{,}40 \times 10^{-6}\ ^\circ\text{C}^{-1}$, respectivamente \cite{serway}.

\label{sec:introduccion} % Let's add label for consistency

La tasa de dilatación lineal depende fundamentalmente de la composición del material y de la intensidad de sus enlaces químicos. Los enlaces atómicos más fuertes, característicos de potenciales interatómicos profundos y simétricos, resisten con mayor eficacia la expansión térmica ante el aumento de energía cinética, lo que resulta en coeficientes de dilatación bajos. Por este motivo, diferentes materiales muestran distinta susceptibilidad a la dilatación, siendo crucial estudiar comparativamente su respuesta térmica. Este comportamiento diferencial tiene un impacto crítico en el diseño de ingeniería de estructuras expuestas a fluctuaciones térmicas \cite{serway}.

El presente experimento tiene por objeto determinar empíricamente el coeficiente de dilatación lineal $\alpha$ para tres materiales distintos: aluminio, cobre y vidrio. El sistema experimental consiste en un dilatómetro clásico donde se introduce vapor de agua a una temperatura de \qty{100}{\celsius} a través de tubos huecos de cada material, elevando su temperatura desde una temperatura ambiental de \qty{26}{\celsius}. Para medir el desplazamiento submilimétrico del tubo al expandirse, se emplea el método de la aguja giratoria, en el cual el extremo libre del tubo descansa sobre una aguja circular de diámetro $D$. Al expandirse el tubo, este se traslada y hace rodar sin deslizar a la aguja un ángulo $\Delta\theta$ sobre un dial graduado, de modo que el alargamiento térmico se calcula geométricamente mediante la relación $\Delta L = D\,\Delta\theta$ \cite{halliday}.

Se anticipa que los coeficientes determinados experimentalmente reflejen el orden de dilatación de los materiales y se aproximen a sus referentes bibliográficos. La precisión de este método angular se evaluará mediante la discrepancia porcentual respecto a los valores de referencia para cada material. Asimismo, los resultados permitirán analizar las principales fuentes de error experimental, tales como las pérdidas térmicas por convección a lo largo del tubo, la fricción del rodamiento o la incertidumbre instrumental del dial de giro y de la medición del diámetro de la aguja.